
We measure curvature subspace alignment between Hessian outlier eigenvectors and minority-group gradients. The method consists of three components: identifying sharp curvature subspaces using Marchenko-Pastur theory, computing minority-group gradients, and measuring alignment.

\subsubsection{Marchenko-Pastur Curvature Subspace Detection}

For a neural network with M parameters trained on N samples where M ≥ N, Marchenko-Pastur random matrix theory predicts that eigenvalues of a random matrix follow a bulk distribution with upper edge λ₊ = (1 + √(N/M))². Eigenvalues exceeding λ₊ represent signal—concentrated curvature from data structure—while eigenvalues within the bulk represent noise.

We compute the top-K eigenvalues of the Hessian using Lanczos iteration. For Waterbirds with N = 11,788 samples and ResNet-50 with M = 25.6M parameters, we compute K = 100 eigenvalues. We fit a Marchenko-Pastur distribution to the empirical spectrum using maximum likelihood estimation to obtain bulk edge λ₊. Eigenvalues λᵢ > λ₊ define the outlier subspace S_out spanned by corresponding eigenvectors vᵢ.

\subsubsection{Minority-Group Gradient Computation}

We compute per-sample gradients g(x) = ∇_θ L(f(x; θ), y) on minority-group samples. For Waterbirds, minority groups consist of landbirds on water backgrounds (184 samples) and waterbirds on land backgrounds (56 samples), totaling 240 samples. We average these gradients: g_minority = (1/|B_minority|) Σ g(x). We compute g_majority similarly for the 4,555 majority-group samples.

Computing this diagnostic metric requires identifying which samples belong to minority groups. While group labels are not required during model training, they are necessary for post-training diagnostic computation. This distinguishes our approach from fully unsupervised methods but avoids the annotation costs of group-labeled training.

\subsubsection{Alignment Metric}

We measure the fraction of minority-gradient variance that projects onto the outlier subspace:

A(θ) = ||P_{S_out} g_minority||² / ||g_minority||²

where P_{S_out} is the projection operator onto the outlier subspace. This metric ranges from 0 to 1, where higher values indicate that minority gradients point into sharp curvature directions. We compute the projection as p = Σᵢ (vᵢᵀ g_minority) vᵢ and calculate alignment as ||p||² / ||g_minority||².
